\section*{第 \ref{ch:b3:relativistic-dynamics} 章——相对论动力学}

\begin{solution}{exo:b3:relativistic-dynamics:1}
电子：$\gamma = 7.09$；$pc = \gamma\beta\,mc^2 = \qty{3.59}{MeV}$，
故 $p = \qty{3.59}{MeV}/c$；$E = \qty{3.62}{MeV}$，$E_k =
\qty{3.11}{MeV}$。质子：$E = 938 + 1000 = \qty{1938}{MeV}$，$\gamma
= 2.07$，$pc = \sqrt{E^2 - (mc^2)^2} = \qty{1696}{MeV}$，$\beta =
pc/E = 0.875$。
\end{solution}

\begin{solution}{exo:b3:relativistic-dynamics:2}
(a) $\num{e-3} \times \num{9e16} = \qty{9e13}{J}$。(b) $\num{6e13}/
\num{9e16} \approx \qty{0.7}{g}$——炸弹转化的
质量不到一克。(c) $L/c^2 = \qty{4.2e9}{kg/s}$；在 \qty{e10}{yr}
（$\num{3.2e17}{}\,\unit{s}$）：$\num{1.3e27}{}\,\unit{kg}$，约占
太阳质量的 $0.07\%$。(d) $\num{5e4}/\num{9e16} \approx
\qty{0.6}{ng}$——真实存在，且永远无法称量。
\end{solution}

\begin{solution}{exo:b3:relativistic-dynamics:3}
(a) 每秒的动量 $p = P/c$，反冲力为 $F = P/c = \qty{1.7e-11}{N}$。
(b) 反射使动量传递加倍：$F = 2\Phi A/c =
2 \times 1400 \times 100/\num{3e8} \approx \qty{0.9}{mN}$。(c) $a
\approx \qty{9e-5}{m/s^2}$；经过 \qty{2.6e6}{s} 后：$v \approx
\qty{240}{m/s}$——起步缓慢，但推进器永不耗尽。(d)
太阳光的动量（连同太阳风）持续将尘埃
推向外侧：彗尾背离太阳伸展，而非
位于彗星之后。
\end{solution}

\begin{solution}{exo:b3:relativistic-dynamics:4}
(a) 由 $m = E_0/c^2$：$\num{0.511e6} \times \num{1.6e-19}/\num{9e16} =
\qty{9.1e-31}{kg}$。(b) $\num{e9} \times \num{1.6e-19}/\num{3e8} =
\qty{5.3e-19}{kg.m/s}$。(c) $\sqrt s = \qty{4}{GeV}$ 为
质心系能量。(d) 取 $c = 1$，能量、动量与
质量使用同一单位，每个公式中的 $c$ 皆被消去；若要回到
国际单位制，需补入使量纲一致的唯一 $c$ 的幂次。
\end{solution}

\begin{solution}{exo:b3:relativistic-dynamics:5}
(a) $2E_k/mc^2$：$1.96$、$3.9$、$17.6$、$59$——第一个之后
皆荒谬。(b) $1 - (1 + E_k/mc^2)^{-2}$：$0.74$、$0.89$、$0.990$、
$0.9989$。(c) $8.4/0.99946c = \qty{28.0}{ns}$；若按牛顿力学，$7.7c$ 将
给出 \qty{3.7}{ns}。(d) 量热计证明电子确实将
全部 \qty{15}{MeV} 带入靶中：能量俱在，
然而速率已不再增长——功如今换来的是动量和能量，
而非速度。
\end{solution}

\begin{solution}{exo:b3:relativistic-dynamics:6}
(a) $m_\pi c^2 = \sqrt{p^2c^2 + m_\mu^2c^4} + pc$；分离根式
并平方，得 $pc = (m_\pi^2 - m_\mu^2)c^4/2m_\pi c^2 =
(139.6^2 - 105.7^2)/(2 \times 139.6) = \qty{29.8}{MeV}$。(b) $E_\mu =
\sqrt{29.8^2 + 105.7^2} = \qty{109.8}{MeV}$：$E_k = \qty{4.1}{MeV}$，
$\beta = 29.8/109.8 = 0.27$。(c) $E_\nu = pc = \qty{29.8}{MeV}$。(d)
作洛伦兹变换：$E_{\text{实验室}} = \gamma(E^* \pm \beta p^*c)$，取 $\beta
\approx 1$：介于 $50 \times (109.8 - 29.8) = \qty{4.0}{GeV}$ 与
$50 \times (109.8 + 29.8) = \qty{7.0}{GeV}$ 之间。
\end{solution}

\begin{solution}{exo:b3:relativistic-dynamics:7}
(a) 对两个无静质量的量子，$M^2c^4 = 2E_1E_2(1 - \cos\theta) =
2 \times 200^2 \times \tfrac12$：$M = \qty{200}{MeV}/c^2$。(b)
恰为 $m_{\pi^0}$——不变质量即粒子的质量，
在任何参考系中都相同。(c) 对事件中每一对光子计算 $M$：
随机配对的结果平缓分布，而粒子的真实衰变产物则
在其质量处堆积，形成峰。(d) $\theta = 0$ 时：$M = 0$；平行光束
没有静止参考系——整体以 $c$ 运动，正如其中
每个光子一样。
\end{solution}

\begin{solution}{exo:b3:relativistic-dynamics:8}
(a) $\gamma = \num{6.8e12}/\num{9.38e8} = 7250$；$c - v \approx
c/2\gamma^2 = \qty{2.9}{m/s}$。(b) $f = v/L \approx \num{3e8}/27000 =
\qty{11.1}{kHz}$：每秒约一万一千圈。(c) $\num{3e14}
\times \qty{6.8}{TeV} = \qty{3.3e8}{J} \approx \qty{78}{kg}$ 的 TNT 当量
——集中于一束发丝般细的光束中。(d) $2 \times \num{3.3e8}/\num{9e16} \approx
\qty{7}{\micro g}$：待转化物质的工作储量。
\end{solution}

\begin{solution}{exo:b3:relativistic-dynamics:9}
(a) 光子的不变质量为 $0$；$e^+e^-$ 对的不变质量至少为
$2m_{\text{e}}c^2$。孤立衰变中不变量不能改变：
禁止。（等价地：任何参考系中动量均不能守恒。）(b)
对头碰撞时，$M^2c^4 = 4E\epsilon$；电子对产生需满足 $Mc^2 \ge
2m_{\text{e}}c^2$，即 $E\epsilon \ge (m_{\text{e}}c^2)^2$。(c)
另需一个 \qty{511}{keV} 光子；若为 \qty{50}{keV}，则伴随光子的能量为
$0.511^2/0.05 = \qty{5.2}{MeV}$。(d) 一个 TeV 光子在
普通星光中的阈值为 $(m_{\text{e}}c^2)^2/E \sim \qty{0.3}{eV}$：天空本身会在宇宙学
距离上吸收 TeV 级 $\gamma$ 射线——
宇宙并非在一切能量下都透明。
\end{solution}

\begin{solution}{exo:b3:relativistic-dynamics:10}
(a) $P_{e'} = P_\gamma + P_e - P_{\gamma'}$；两边平方
（取不变量）：$m_{\text{e}}^2c^4 = m_{\text{e}}^2c^4 + 2P_\gamma\!
\cdot\!P_e - 2P_{\gamma'}\!\cdot\!P_e - 2P_\gamma\!\cdot\!P_{\gamma'}$。
(b) 由 $P_\gamma\!\cdot\!P_e = E m_{\text{e}}$，
$P_{\gamma'}\!\cdot\!P_e = E'm_{\text{e}}$ 及 $P_\gamma\!\cdot\!
P_{\gamma'} = EE'(1 - \cos\theta)/c^2$，
得 $m_{\text{e}}c^2(E - E') = EE'(1 - \cos\theta)$；以波长表示
（$E = hc/\lambda$）即为所述的波长偏移。(c) $h/m_{\text{e}}c =
\qty{2.43}{pm}$，至多 \qty{4.9}{pm}：对可见光，约为
其 $10^5$ 分之一（不可见），对 \qty{100}{pm} 的 X 射线则约占百分之几
（可测）。(d) 光子像台球一样碰撞——
携带动量 $h/\lambda$，并在\emph{单个}事件中交换，
而不只是以团块形式交换能量。
\end{solution}

\begin{solution}{exo:b3:relativistic-dynamics:11}
(a) 对头碰撞时 $M^2c^4 = m_{\text{p}}^2c^4 + 4E\epsilon$（质子为
超相对论性）；阈值在 $M = m_\Delta$ 处。(b) $E =
(1232^2 - 938^2)\,\unit{MeV^2}/(4 \times \num{6e-4}\,\unit{eV})
\approx \qty{2.7e20}{eV}$（对平均光子而言）——宇宙微波背景的热尾
将有效截断降至
$\sim\qty{6e19}{eV}$。(c) 高于截断的质子将能量
损失给 $\Delta$，在 $\sim\qty{e8}{ly}$ 内：对遥远源，能谱应在
\qty{6e19}{eV} 附近截止——即观测到的 GZK 压低。
(d) 其源必须既是非凡的加速器，
又在宇宙尺度上\emph{较近}——位于我们的超星系团之内，
这也是其起源至今仍在探寻的部分原因。
\end{solution}

\begin{solution}{exo:b3:relativistic-dynamics:12}
(a) 能量：$m_{\text{i}}c^2 = \gamma m_{\text{f}}c^2 + E_\ell$；
动量：$\gamma m_{\text{f}}\beta c = E_\ell/c$。消去
$E_\ell$：$m_{\text{i}} = \gamma(1 + \beta)m_{\text{f}} =
m_{\text{f}}\sqrt{(1+\beta)/(1-\beta)}$。(b) $\sqrt{19} = 4.4$；
加速\emph{和}制动使之平方：$19$。(c) 消耗的
质量 $m_{\text{i}} - m_{\text{f}} = 3.4\,m_{\text{f}}$ 必须是
被湮灭的燃料，其中一半为反物质：\qty{1.7}{kg} 反物质
每交付一千克（单程，无制动）。(d) 以全球年产量的纳克级计，
光子火箭仍只是算术，
而非工程。
\end{solution}

\begin{solution}{pb:b3:relativistic-dynamics:1}
\textbf{1.} $E^2 - p^2c^2 = \gamma^2m^2c^4(1 - \beta^2) = m^2c^4$。
\textbf{2.} $E_{\text{总}}$ 与 $\vect p_{\text{总}}$ 的变换方式
与单个粒子的 $E$ 和 $\vect p$ 相同（它们是求和量），因此同一
代数给出不变量；在 $\vect p_{\text{总}} = \vect 0$ 处，它
即总能量的平方：$Mc^2 = E_{\text{CM}}$。
\textbf{3.} $(E + m_{\text{p}}c^2)^2 - p^2c^2 = m_{\text{p}}^2c^4 +
2Em_{\text{p}}c^2 + m_{\text{p}}^2c^4$。
\textbf{4.} $E \to m_{\text{p}}c^2$ 给出 $M = 2m_{\text{p}}$；
$E \gg m_{\text{p}}c^2$ 时，$Mc^2 \to \sqrt{2Em_{\text{p}}c^2}$。
\textbf{5.} 产物的任何相对运动，都会在质心系中
在其静质量之上再增加动能；最小的 $Mc^2$——
即阈值——使它们全部处于相对静止。
\textbf{6.} 实验室系中的动量 $\vect p_{\text{总}} \neq \vect 0$ 必须
在碰撞中保持：产物注定要运动，其
动能被锁住，无法用于产生质量。
\textbf{7.} $p + p \to p + p + p + \bar p$：电荷 $+2 \to +2$；
重子数 $+2 \to 1 + 1 + 1 - 1 = +2$。单独产生 $\bar p$，或
伴随少于一个完整重子的粒子产生，都会破坏其中一条守恒。
\textbf{8.} 四个质子（相对静止）：$Mc^2 =
4m_{\text{p}}c^2$。
\textbf{9.} $16m_{\text{p}}^2c^4 = 2m_{\text{p}}^2c^4 +
2Em_{\text{p}}c^2$：$E = 7m_{\text{p}}c^2$，$T = 6m_{\text{p}}c^2 =
\qty{5.63}{GeV}$。
\textbf{10.} 新的静质量：$2m_{\text{p}}c^2$，占 $T =
6m_{\text{p}}c^2$ 的三分之一；其余三分之二以
产物动能的形式飞出，并受动量守恒保护。
\textbf{11.} 正面碰撞时，$Mc^2 = 2(T' + m_{\text{p}}c^2) =
4m_{\text{p}}c^2$：$T' = m_{\text{p}}c^2 = \qty{0.94}{GeV}$，每束
——比固定靶束流低六倍，
总动能低十二倍。
\textbf{12.} 1950 年代的束流过于稀薄，难以彼此有效碰撞；
固体靶每立方厘米可提供 $10^{22}$ 个质子。
代价是 $\qty{5.6}{GeV}$，而非
$2 \times \qty{0.94}{GeV}$。
\textbf{13.} $E = 6.2 + 0.94 = \qty{7.14}{GeV}$；$pc =
\sqrt{7.14^2 - 0.938^2} = \qty{7.08}{GeV}$。
\textbf{14.} $p = \num{7.08e9} \times \num{1.6e-19}/\num{3e8} =
\qty{3.8e-18}{kg.m/s}$：$r = p/qB = \qty{15.1}{m}$——与
实际建成的机器一致。
\textbf{15.} $\beta = pc/E = 0.991$；轨道 $2\pi r = \qty{95}{m}$：
$f = \qty{3.1}{MHz}$。
\textbf{16.} 在 \qty{10}{MeV} 时，$\beta = 0.145$：$f =
\qty{0.46}{MHz}$。随着能量升高，速率（从而 $f$）与
动量（从而固定半径处所需的场）都在变化：场
与射频频率必须同步升高——这正是同步
加速器的定义性技巧。
\textbf{17.} $\num{6.19e9}/1500 \approx \num{4e6}$ 圈；在
平均兆赫量级频率下，一个加速周期的量级
为一至两秒。
\textbf{18.} 反质子阈值 $T = 6m_{\text{p}}c^2 =
\qty{5.6}{GeV}$，加上余量后：本题的核算
字面上正是这台机器的能量乃至其名称所为之设计的。
\textbf{19.} $r = p/qB$ 中不含质量：磁铁使相等的动量
获得同样的偏转，无论是何种粒子——因此所选束流仍
混有 π 介子、K 介子及（少量）反质子，
三者动量均为 $\qty{1.19}{GeV}/c$。
\textbf{20.} $\bar p$：$E = \sqrt{1.19^2 + 0.938^2} =
\qty{1.51}{GeV}$，$\beta = 0.79$。$\pi^-$：$E = \sqrt{1.19^2 +
0.140^2} = \qty{1.20}{GeV}$，$\beta = 0.993$。
\textbf{21.} $t_{\bar p} = 12/(0.785c) = \qty{51}{ns}$；$t_\pi =
\qty{40}{ns}$：\qty{11}{ns} 的间隔要求纳秒级计时——
1955 年勉强可以做到。
\textbf{22.} 计数器被设定为对快速 π 介子触发，而对慢速
反质子保持静默（速度否决）；在 $1$ 比 $44000$ 的
“大海捞针”中，只有两个独立特征（飞行时间\emph{与}
切伦科夫阈值）才能排除巧合假信号。
\textbf{23.} 刚高于阈值时，产物几乎在
相对静止时产生：反应几乎没有相空间，而 π 介子
产生远高于\emph{其}阈值，十分丰富——稀有性
由决定代价的同一运动学所保证。
\textbf{24.} $2m_{\text{p}}c^2 = \qty{1.9}{GeV}$，每次湮灭，
通常生成少量 π 介子，这些介子继续衰变为光子、μ 子
与中微子。
\textbf{25.} 四维动量守恒将反质子在固定靶上的代价定为
$6m_{\text{p}}c^2 = \qty{5.6}{GeV}$；
\qty{6.2}{GeV}、\qty{15.1}{m} 的 Bevatron 支付了这笔代价；飞行时间与
切伦科夫否决在每 $44000$ 个冒充者中找出一个
反世界粒子——1959 年诺贝尔奖确认了这份账目。
\end{solution}
